\subsection{Resumen de las interacciones}

\begin{center}
\footnotesize
\setlength{\tabcolsep}{4pt}
\renewcommand{\arraystretch}{1.3}
\begin{tabularx}{\textwidth}{@{}c l >{\raggedright\arraybackslash}X >{\raggedright\arraybackslash}p{1.9cm} >{\raggedright\arraybackslash}p{1.6cm} >{\raggedright\arraybackslash}p{2.5cm} >{\raggedright\arraybackslash}p{2.2cm}@{}}
\toprule
\textbf{N.º} & \textbf{De $\to$ A} & \textbf{Mensaje (petición / respuesta)} & \textbf{Mecanismo} & \textbf{Tipo} & \textbf{Protocolo} & \textbf{Middleware} \\
\midrule
1 & P1 $\to$ P2 & \texttt{BUSCAR\_VUELOS}(origen, destino, fecha) / lista de vuelos & Sockets & Síncrona & TCP, puerto 5000 & Ninguno (socket de Python, JSON) \\
2 & P2 $\to$ P4 & Consulta de vuelos y asientos libres / filas encontradas & Sockets (driver de BD) & Síncrona & TCP, puerto 5432, protocolo de PostgreSQL & psycopg \\
3 & P1 $\to$ P2 & \texttt{RESERVAR}(vuelo, asiento, pasajero, id\_solicitud) / código de reserva o rechazo & Sockets & Síncrona & TCP, puerto 5000 & Ninguno (socket de Python, JSON) \\
4 & P2 $\to$ P4 & Transacción: bloquear el asiento y guardar la reserva como PENDIENTE (10 min) / OK & Sockets (driver de BD) & Síncrona & TCP, puerto 5432, protocolo de PostgreSQL & psycopg \\
5 & P2 $\to$ P3 & \texttt{ProcesarPago}(reserva, monto, token, clave) / APROBADO o RECHAZADO & RPC & Síncrona, deadline de 3 s & HTTP/2 sobre TCP, puerto 50051 & gRPC y Protocol Buffers \\
6 & P2 $\to$ P4 & Cambiar la reserva a CONFIRMADA o CANCELADA / OK & Sockets (driver de BD) & Síncrona & TCP, puerto 5432, protocolo de PostgreSQL & psycopg \\
7 & P2 $\to$ P1 & \texttt{NOTIFICAR}(cambio de vuelo o reserva vencida) / sin respuesta & Sockets & Asíncrona & TCP, puerto 5000 & Ninguno (socket de Python, JSON) \\
\bottomrule
\end{tabularx}
\end{center}

Las interacciones 1, 3 y 7 usan sockets y la 5 usa RPC con gRPC, que son los dos mecanismos que vamos a programar y comparar en PE-U1.

\subsection{Justificación de cada interacción}

\begin{description}
  \item[1. Buscar vuelos.] El cliente abre una conexión TCP con el servidor y la deja abierta mientras usa la app. Elegimos sockets porque con la conexión ya abierta cada búsqueda no tiene que volver a conectarse, y así la respuesta llega más rápido. Es síncrona porque el usuario no puede seguir hasta ver los vuelos. Los mensajes son pequeños (un JSON por línea), así que aunque haya muchas búsquedas el volumen no es un problema. TCP asegura que el mensaje llegue completo y en orden.

  \item[2. Consultar vuelos.] El servidor le pide a PostgreSQL los vuelos y los asientos libres. Por debajo también es un socket TCP, pero no lo programamos nosotros: lo maneja el driver psycopg con el protocolo propio de PostgreSQL. Es síncrona porque el servidor necesita esos datos para contestar la interacción 1.

  \item[3. Reservar.] Va por la misma conexión que la búsqueda. Es síncrona porque el usuario espera saber si su reserva quedó o no. Cada solicitud lleva un \texttt{id\_solicitud}: si la conexión se corta y el cliente vuelve a enviar lo mismo, el servidor reconoce el id y no crea una segunda reserva.

  \item[4. Retener el asiento.] Aquí lo que más importa es la consistencia: si dos personas piden el mismo asiento al mismo tiempo, solo una lo puede obtener. Por eso el servidor bloquea la fila del asiento dentro de una transacción y guarda la reserva como PENDIENTE por 10 minutos. La transacción se cierra enseguida, antes de cobrar, para que el asiento no quede bloqueado mientras se espera al servicio de pagos.

  \item[5. Procesar el pago.] Es la interacción con RPC. Usamos gRPC porque es una comunicación interna entre dos servicios del sistema: el archivo \texttt{.proto} define exactamente qué datos se envían y qué se recibe, Protocol Buffers manda los datos en binario y HTTP/2 reutiliza la conexión, por eso la latencia es baja. Es síncrona porque sin el resultado del pago no se puede confirmar la reserva. Le ponemos un deadline de 3 segundos para que el servidor no se quede esperando si pagos no contesta, y una clave de idempotencia para poder reintentar una vez sin cobrar dos veces. Al servicio de pagos solo le llega un token de la tarjeta, nunca el número.

  \item[6. Confirmar o liberar.] Con la respuesta de pagos, el servidor cambia la reserva a CONFIRMADA, o a CANCELADA y libera el asiento. Así lo que queda guardado en la base de datos coincide con lo que pasó con el pago. Recién ahí el servidor contesta al cliente la interacción 3.

  \item[7. Notificar.] El servidor avisa al cliente cuando cambia algo de su vuelo, como un retraso o un cambio de puerta, o cuando vence una reserva pendiente. Es asíncrona: el servidor manda el aviso por la conexión abierta y sigue trabajando sin esperar respuesta. Si el cliente está desconectado el aviso se pierde, pero no es grave porque al volver a conectarse consulta el estado de su reserva. Si estos avisos fueran críticos usaríamos una cola de mensajes como RabbitMQ, que guarda el mensaje hasta que se entregue.
\end{description}
